% =====================================================================
% Section 4 -- Evaluation
% Usage: % =====================================================================
% Section 4 -- Evaluation
% Usage: % =====================================================================
% Section 4 -- Evaluation
% Usage: % =====================================================================
% Section 4 -- Evaluation
% Usage: \input{sec4_evaluation} after sec3_system (shares its macros).
% Bibliography keys used below (add to your .bib):
%   menegazzo2021pvs, min2023factscore, wiseman2017challenges, kent1964wep,
%   mielke2022linguistic, williams1949designs, judd2012stimuli, jian2000trust,
%   christensen_ordinal, wobbrock2011art, braun2006thematic, demontjoye2013unique,
%   thompson2026localising, vandeemter2005templates
% =====================================================================
\providecommand{\R}{\mathbb{R}}
\providecommand{\ind}{\mathbf{1}}
\providecommand{\TBD}[1][]{\textcolor{red}{[TBD#1]}}

\section{Evaluation}
\label{sec:eval}

We evaluate RoadScribe on five questions.
\begin{description}
\item[RQ1] Does Layer~1 recover vehicle-frame vibration and road events from a phone, and are its outputs grounded in vibration rather than in speed?
\item[RQ2] Does map-informed segmentation recover road-surface boundaries better than signal-only or map-only segmentation?
\item[RQ3] How faithful are the generated narratives to the ledger and to ground truth, and how do errors split between sensing and language?
\item[RQ4] Where does a grounded SLM add coverage over a template narrator, and at what fidelity cost?
\item[RQ5] How do people judge the narratives for accuracy, usefulness, trust and annoyance, and do they notice errors?
\end{description}

% ---------------------------------------------------------------------
\subsection{Data}
\label{sec:data}

\paragraph{The PVS dataset.}
The Passive Vehicular Sensors dataset~\cite{menegazzo2021pvs} contains nine trips that cover three routes (13.8, 11.6 and 10.7~km), each driven once by three vehicle--driver pairs.
In total there are 180.1~min and 108.6~km of driving.
IMU modules at the dashboard, above the suspension and below the suspension log at 100~Hz on each side of the vehicle, GNSS logs at 1~Hz, and an environment camera records each trip.
Labels mark the surface type (asphalt, cobblestone, dirt; 54 transitions) and speed bumps (60 traversals of 20 physical bumps).
We use the left-side logger: the dashboard unit serves as the phone proxy, and the suspension units serve as reference sensors that a phone does not have.
The left and right loggers are not synchronized.
Cross-correlating their yaw rates shows offsets of up to 0.54~s that drift within a trip, so we do not average the two sides.
PVS uses fixed modules, which exercise neither phone handling nor holder dynamics.

\paragraph{Phone corpus.}
To cover what PVS cannot, we recorded \TBD{} drives (\TBD{}~km, \TBD{} vehicles) with two phones.
One phone was in a dashboard holder and recorded video; the other was fixed to the seat rail as a floor reference.
Both logged at \TBD{}~Hz and were synchronized by tap events and yaw-rate cross-correlation (residual offset \TBD{}~ms).
The corpus includes \TBD{} repetitions of one route by one driver and \TBD{} drives with deliberate phone handling.
Two annotators marked speed breakers, potholes, rough patches, surface changes and rail crossings from video without seeing any narrative (Cohen's $\kappa=\TBD{}$ on \TBD{}\% double-coded).

\paragraph{Splits.}
All learned components, thresholds and prompts are fit with routes held out.
On PVS we use leave-one-route-out folds, and for vehicle generalization we test each (route, vehicle) trip using only the four trips that share neither its route nor its vehicle.
Evaluation questions, test personas and capsules from test drives never appear in prompts, exemplars or training data.

% ---------------------------------------------------------------------
\subsection{Narrators Compared}
\label{sec:conditions}

\begin{table}[t]
\caption{Narrators compared. All share Layers 1--2 and the ledger.}
\label{tab:conditions}
\begin{tabular}{lp{0.62\linewidth}}
\toprule
Code & Description \\
\midrule
T & Template NLG pipeline: rule-based content selection, aggregation, referring expressions, ISO-class lexical choice, history comparison, a fallback for every missing field (\TBD{} rules). \\
T+V & T with sampled lexical and ordering variation (controls for variety alone). \\
T$\rightarrow$P & T followed by SLM paraphrase, as in deployed driver-feedback systems~\cite{thompson2026localising}. \\
S-free & SLM given the capsule, free generation, no plan or verifier. \\
S-PGR & Plan--ground--realize with verification and template fallback, on device (ours). \\
C-PGR & S-PGR protocol with a cloud model (\TBD{}), as an upper reference. \\
\bottomrule
\end{tabular}
\end{table}

% ---------------------------------------------------------------------
\subsection{Metrics}
\label{sec:metrics}

\paragraph{Claims.}
A narrative $N$ is decomposed into atomic claims $\mathcal X(N)$ of type severity, numeric, event, anchor, comparison or hedge, in the spirit of~\cite{min2023factscore,wiseman2017challenges}.
The same LLM-based extractor is applied to every narrator.
On \TBD{} claims stratified by narrator it agreed with two annotators at $\kappa=\TBD{}$.
A claim is supported by reference $R$ when its value, unit, place and direction agree with $R$ within type-specific tolerances.
We use two references: the ledger $L$, which isolates the language layer, and ground truth $G$, which measures the system end to end:
\begin{equation}
P_R(N)=\frac{1}{|\mathcal X(N)|}\sum_{c\in\mathcal X(N)}\mathrm{supp}_R(c),\qquad R\in\{L,G\}.
\end{equation}
We report claims about facts marked \textsc{missing} separately; their rate should be zero.

\paragraph{Kinetic--semantic fidelity.}
Let $\mathcal S$ be the severity claims of a narrator pooled over the corpus.
Each claim $c\in\mathcal S$ pairs a measured intensity $x_c$ ($a_v$ for how the ride felt, $\tilde a_{w,z}$ for the road) with an expressed class $\hat k_c\in\{0,\dots,5\}$ from a fixed lexicon of ISO comfort descriptors.
With $R_c$ and $Q_c$ the mid-ranks of $x_c$ and $\hat k_c$,
\begin{equation}
\rho_{\mathrm{KSF}}=\frac{\sum_{c}(R_c-\bar R)(Q_c-\bar Q)}{\sqrt{\sum_{c}(R_c-\bar R)^2\sum_{c}(Q_c-\bar Q)^2}} .
\end{equation}
Because $\rho_{\mathrm{KSF}}$ ignores monotone distortion, we also report band consistency and an overstatement index,
\begin{equation}
\mathrm{BC}=\frac1{|\mathcal S|}\sum_{c\in\mathcal S}\ind\bigl[x_c\in\mathcal B(\hat k_c)\bigr],\qquad
\mathrm{OI}=\frac1{|\mathcal S|}\sum_{c\in\mathcal S}\bigl(\hat k_c-k^\star(x_c)\bigr),\qquad
k^\star(x)=\sum_{j=1}^5\ind[x\ge\theta_j],
\end{equation}
where $\mathcal B(k)$ is the (overlapping) ISO band of class $k$ and $\boldsymbol\theta$ are the cut-points of Section~\ref{sec:layer1}.

\paragraph{Anomaly grounding.}
Narrated event mentions $\hat{\mathcal E}$ are expanded by count; vague quantifiers expand to the lower bound of a fixed lexicon.
Each mention has a type and a resolved along-track interval $I(\hat e)$, and each reference event $e\in\mathcal E^R$ has a type and a position $s(e)$.
A mention and a reference event may match when their types agree and $\mathrm{dist}(s(e),I(\hat e))\le\delta_s$.
With $M^\star$ a maximum bipartite matching,
\begin{equation}
\mathrm{AGP}_R=\frac{|M^\star|}{|\hat{\mathcal E}|},\qquad
\mathrm{AGR}_R=\frac{|M^\star|}{|\mathcal E^R|},\qquad
\mathrm{wAGR}_R=\frac{\sum_{e\in M^\star}\mathrm{MTVV}_e}{\sum_{e\in\mathcal E^R}\mathrm{MTVV}_e}.
\end{equation}
For $R=L$, recall is computed over the must-mention events.
For $R=G$, the references are the labelled bumps (PVS) and the annotated events (phone corpus).
We use $\delta_s=25$~m, which covers label timing uncertainty and map-match error, and report 10 and 50~m in the supplement.
Crossing the two references attributes every error to perception (missed or false detection), selection (omitted) or verbalization (unsupported).

\paragraph{Calibration of hedges.}
For event claims with hedge $h\in\{\text{assertive},\text{likely},\text{possibly}\}$, let $P_G(h)$ be their precision against ground truth and $p_h$ the probability participants associate with the hedge, elicited in a short survey (cf.~\cite{kent1964wep,mielke2022linguistic}). Then
\begin{equation}
\mathrm{ECE}_{\mathrm{ling}}=\sum_h\frac{n_h}{n}\,\bigl|P_G(h)-p_h\bigr| .
\end{equation}

\paragraph{Novelty--change alignment.}
For consecutive drives $(d,d')$ of one route by one driver, let $\Delta_{\mathrm{kin}}$ be the mean absolute change of $\log a_v$ over 100~m bins plus the event-set Jaccard distance.
Let $\Delta_{\mathrm{claim}}$ be the Jaccard distance between canonicalized claim sets.
Then $\mathrm{NCA}=\mathrm{Spearman}(\Delta_{\mathrm{kin}},\Delta_{\mathrm{claim}})$ over pairs.
Lexical variation independent of content, $\mathrm{LV}=1-\mathrm{BLEU}$ between narratives with numbers and names masked, is reported alongside; variety that does not track change raises LV but not NCA.

\paragraph{Question answering and cost.}
For questions we report the correct-answer rate (judged against the ledger by raters blind to the narrator), the unsupported-claim rate, and abstention precision and recall on unanswerable questions.
System cost comprises time to first token, decode rate, end-to-end latency (p50/p95), peak memory, energy per narrative, plan validity, verifier rejection rate and template fallback rate.

% ---------------------------------------------------------------------
\subsection{Layer 1 and Layer 2}
\label{sec:l1eval}

\paragraph{Orientation.}
On PVS, the heading estimate of Eq.~\eqref{eq:yaw} recovered the known mount: $87.9^\circ$--$92.5^\circ$ for two vehicles (nominal $90^\circ$).
For the third vehicle it gave $103.1^\circ$--$105.5^\circ$ consistently across its three routes, indicating a mounting offset of about $14^\circ$.
Under 20 random three-dimensional rotations of the input per trip, the estimate rotated with the data to within $0.63^\circ$.
On the phone corpus, handling detection and re-estimation achieve \TBD{}.

\paragraph{Do maneuvers inflate vertical vibration?}
Table~\ref{tab:maneuver} compares two-second windows on asphalt (no speed bump within $\pm2$~s, 5--20~m\,s$^{-1}$) across maneuver states, controlling for speed and trip.
Braking lowered dashboard vertical vibration and turning had no detectable effect.
Acceleration added high-band energy at the dashboard that the below-suspension sensor did not see, consistent with powertrain vibration.
This is why the gate of Eq.~\eqref{eq:gate} targets acceleration and stop-and-go events and never vetoes on braking.
A wheelbase-lag cue (front then rear axle) did not separate road events: the lag between the two largest impacts did not scale with $1/v$ (Spearman $|\rho|\le0.12$, $n=60$), so it is not used.

\begin{table}[t]
\caption{Difference in log RMS vertical acceleration relative to cruising (PVS, asphalt, 2-s windows; OLS with log speed and trip effects, trip-clustered standard errors). Cruise $n=225$.}
\label{tab:maneuver}
\begin{tabular}{lcccc}
\toprule
Maneuver ($n$) & Dash 0.5--4 Hz & Dash 4--25 Hz & Wheel 4--25 Hz & Dash $-$ wheel, 4--25 Hz\\
\midrule
Brake (75) & $-0.18^{***}$ & $-0.13^{*}$ & $-0.16$ & $+0.03$\\
Accelerate (23) & $-0.18$ & $-0.06$ & $-0.37^{**}$ & $+0.31^{***}$\\
Turn (92) & $+0.07$ & $-0.04$ & $+0.06$ & $-0.10^{***}$\\
\bottomrule
\multicolumn{5}{l}{\footnotesize $^{*}p<.05$, $^{**}p<.01$, $^{***}p<.001$. Speed elasticity $\hat\gamma$: 0.32, 0.36, 0.62 (columns 1--3).}
\end{tabular}
\end{table}

\paragraph{Events and segmentation.}
Bump detection on held-out routes reaches recall \TBD{} at \TBD{} false alarms per km.
For comparison, a speed-only detector reaches \TBD{} and $\mathrm{RMS}(a_z)/\sqrt v$ reaches \TBD{}.
The gate changes false alarms by \TBD{} at fixed recall.
Against the 54 surface transitions, MI-PELT reaches boundary F1 \TBD{} at 25~m, compared with \TBD{} for constant-penalty PELT, \TBD{} for map-only boundaries and \TBD{} for PELT on speed alone.
OpenStreetMap carries a surface tag on \TBD{}\% of the matched length and a \texttt{traffic\_calming} node at \TBD{} of the 20 bumps.
Results on PVS cover three route clusters and are descriptive; inferential comparisons use the phone corpus.

% ---------------------------------------------------------------------
\subsection{Narrative Benchmark}
\label{sec:bench}

\paragraph{Tasks.}
Five task families probe where a template narrator suffices and where it does not.
\textbf{E1} answers questions from a sealed half of \TBD{} questions elicited from drivers before the system was built; the template baseline uses an intent classifier over intents designed on the other half.
\textbf{E2} writes for audiences that no template or exemplar targets (road engineer, motion-sensitive passenger, fleet manager); the template baseline uses the nearest designed audience.
\textbf{E3} narrates degraded or contradictory capsules: no gyroscope, a 60~s GNSS gap, failed map matching, map--signal disagreement, low-confidence events.
\textbf{E4} narrates repeated drives of one route.
\textbf{E5} grows the template narrator rule by rule on the development split and traces its coverage against the number of rules.
Thresholds that would count against our hypotheses were fixed before the sealed sets were used (\TBD[: registration statement]).

\paragraph{Results.}
\TBD[: Table with $P_L$, $P_G$, hallucinated-missing rate, $\rho_{\mathrm{KSF}}$, BC, OI, AGP/AGR ($L$ and $G$), $\mathrm{ECE}_{\mathrm{ling}}$, NCA, LV, QA correctness, fallback rate per narrator; frontier plot of $P_L$ against task coverage; error attribution table.]

% ---------------------------------------------------------------------
\subsection{On-Device Performance}
\label{sec:device}
\TBD[: two phones (one mid-range); model, quantization, time to first token, decode rate, end-to-end latency p50/p95, peak memory, energy per narrative, battery drain of 100--400~Hz logging, throttling over 10 consecutive generations. State which layers run on the phone.]

% ---------------------------------------------------------------------
\subsection{User Study}
\label{sec:study}

\paragraph{Design.}
Twenty-eight licensed drivers each judged narratives from four narrators (T, T+V, T$\rightarrow$P, S-PGR) in a within-subject design.
Narrator order followed a Williams design~\cite{williams1949designs} with sequences (T, T+V, S-PGR, T$\rightarrow$P), (T+V, T$\rightarrow$P, T, S-PGR), (T$\rightarrow$P, S-PGR, T+V, T) and (S-PGR, T, T$\rightarrow$P, T+V), each used by seven participants, which balances first-order carryover.
Every participant saw eight video excerpts of 2--3~min, two per narrator.
A second Latin square rotated excerpts across narrators, so each excerpt appeared with each narrator equally often and excerpts could be modelled as random effects~\cite{judd2012stimuli}.
In two of each participant's eight narratives one fact was altered (a count, a place or a severity), balanced across narrators, to measure whether errors were noticed.
Participants were debriefed about the altered narratives.

\paragraph{Participants and ethics.}
\TBD[: recruitment, inclusion criteria (licensed, $\ge$1 year of regular driving), demographics, compensation, ethics approval number]. Location traces are identifying~\cite{demontjoye2013unique}; they were stored encrypted and are released only truncated by 500~m at each end. No narrative was presented while a participant was driving.

\paragraph{Procedure.}
After consent and a practice trial, participants completed four blocks, one per narrator.
For each excerpt they watched the video, read the narrative and answered the items of Table~\ref{tab:items}.
They then answered whether anything seemed wrong, and if so, what.
Each block ended with the trust checklist of Jian et al.~\cite{jian2000trust}.
In a question-answering task, participants typed three questions about each of two further excerpts and rated the answers, which came from T with an intent classifier or from S-PGR, counterbalanced.
The session closed with a ranking of the narrators and a semi-structured interview, analysed thematically~\cite{braun2006thematic}.

\begin{table}[t]
\caption{Per-narrative items (7-point scale, 1 = strongly disagree, 7 = strongly agree).}
\label{tab:items}
\begin{tabular}{lp{0.78\linewidth}}
\toprule
Code & Item \\
\midrule
ACC1 & The summary accurately described the road conditions in the clip. \\
ACC2 & The places the summary mentioned matched where things happened in the clip. \\
SPEC & The summary was specific enough that I could find the places it mentioned again. \\
USE & This summary would help me decide how to drive this route next time. \\
CLAR & The summary was easy to understand on the first reading. \\
CERT & When the summary was unsure about something, it said so in a way that felt right. \\
NOVEL & The summary told me something about the drive that I had not noticed myself. \\
ANNOY & I would find getting a summary like this after every drive annoying. (reversed) \\
FIT & This summary was written for someone like me. \\
\bottomrule
\end{tabular}
\end{table}

\paragraph{Hypotheses.}
Four hypotheses were fixed before data collection.
H1: claim precision of S-PGR is non-inferior to T, $P_L(\text{S-PGR})-P_L(\text{T})>-0.02$.
H2: S-PGR answers the sealed questions correctly more often than T with an intent classifier.
H3: S-PGR is rated more useful (USE) than T+V.
H4: perceived accuracy, the mean of ACC1 and ACC2, of S-PGR is non-inferior to T within 0.5 scale points.
All other comparisons are exploratory.

\paragraph{Analysis.}
Ratings were analysed with cumulative link mixed models~\cite{christensen_ordinal},
\begin{equation}
\operatorname{logit}P(Y_{ij}\le k)=\theta_k-\bigl(\beta_{\mathrm{cond}(ij)}+\beta_{\mathrm{period}(ij)}+u_i+w_j\bigr),\qquad u_i\sim\mathcal N(0,\sigma_u^2),\ w_j\sim\mathcal N(0,\sigma_w^2),
\end{equation}
for participant $i$ and excerpt $j$, with Holm correction across H2 and H3 and a first-order carryover term as a check.
We also ran aligned-rank-transform ANOVA~\cite{wobbrock2011art} and Wilcoxon signed-rank tests on participant means as robustness checks.
H4 was tested on participant means of unaltered trials and holds if the lower bound of the 90\% confidence interval of the paired difference exceeds $-0.5$.
H1 used a cluster bootstrap over drives.
H2 and error detection used logistic mixed models with random intercepts for question or participant and for excerpt.

\paragraph{Power.}
With 28 participants, a paired comparison has 80\% power for standardized effects $d_z\ge0.55$ at $\alpha=.05$ ($d_z\ge0.65$ at the Holm-adjusted $\alpha=.05/3$).
The non-inferiority test of H4 has 82.5\% power when the margin equals half the standard deviation of paired differences.
The pilot ($n=\TBD{}$) estimated that standard deviation at \TBD{} scale points.

\paragraph{Field deployment.}
\TBD{} participants used RoadScribe on their own phones for 14 days.
Narrators alternated by day between T+V and S-PGR on a randomized balanced schedule.
Narratives were delivered only after the vehicle had been stationary for two minutes.
We logged opens, time to open, dwell time and questions asked, and collected three daily ratings (USE, ANNOY, NOVEL).
We model habituation as the narrator-by-day interaction in a mixed model with random slopes; with this sample size the analysis is descriptive.
 after sec3_system (shares its macros).
% Bibliography keys used below (add to your .bib):
%   menegazzo2021pvs, min2023factscore, wiseman2017challenges, kent1964wep,
%   mielke2022linguistic, williams1949designs, judd2012stimuli, jian2000trust,
%   christensen_ordinal, wobbrock2011art, braun2006thematic, demontjoye2013unique,
%   thompson2026localising, vandeemter2005templates
% =====================================================================
\providecommand{\R}{\mathbb{R}}
\providecommand{\ind}{\mathbf{1}}
\providecommand{\TBD}[1][]{\textcolor{red}{[TBD#1]}}

\section{Evaluation}
\label{sec:eval}

We evaluate RoadScribe on five questions.
\begin{description}
\item[RQ1] Does Layer~1 recover vehicle-frame vibration and road events from a phone, and are its outputs grounded in vibration rather than in speed?
\item[RQ2] Does map-informed segmentation recover road-surface boundaries better than signal-only or map-only segmentation?
\item[RQ3] How faithful are the generated narratives to the ledger and to ground truth, and how do errors split between sensing and language?
\item[RQ4] Where does a grounded SLM add coverage over a template narrator, and at what fidelity cost?
\item[RQ5] How do people judge the narratives for accuracy, usefulness, trust and annoyance, and do they notice errors?
\end{description}

% ---------------------------------------------------------------------
\subsection{Data}
\label{sec:data}

\paragraph{The PVS dataset.}
The Passive Vehicular Sensors dataset~\cite{menegazzo2021pvs} contains nine trips that cover three routes (13.8, 11.6 and 10.7~km), each driven once by three vehicle--driver pairs.
In total there are 180.1~min and 108.6~km of driving.
IMU modules at the dashboard, above the suspension and below the suspension log at 100~Hz on each side of the vehicle, GNSS logs at 1~Hz, and an environment camera records each trip.
Labels mark the surface type (asphalt, cobblestone, dirt; 54 transitions) and speed bumps (60 traversals of 20 physical bumps).
We use the left-side logger: the dashboard unit serves as the phone proxy, and the suspension units serve as reference sensors that a phone does not have.
The left and right loggers are not synchronized.
Cross-correlating their yaw rates shows offsets of up to 0.54~s that drift within a trip, so we do not average the two sides.
PVS uses fixed modules, which exercise neither phone handling nor holder dynamics.

\paragraph{Phone corpus.}
To cover what PVS cannot, we recorded \TBD{} drives (\TBD{}~km, \TBD{} vehicles) with two phones.
One phone was in a dashboard holder and recorded video; the other was fixed to the seat rail as a floor reference.
Both logged at \TBD{}~Hz and were synchronized by tap events and yaw-rate cross-correlation (residual offset \TBD{}~ms).
The corpus includes \TBD{} repetitions of one route by one driver and \TBD{} drives with deliberate phone handling.
Two annotators marked speed breakers, potholes, rough patches, surface changes and rail crossings from video without seeing any narrative (Cohen's $\kappa=\TBD{}$ on \TBD{}\% double-coded).

\paragraph{Splits.}
All learned components, thresholds and prompts are fit with routes held out.
On PVS we use leave-one-route-out folds, and for vehicle generalization we test each (route, vehicle) trip using only the four trips that share neither its route nor its vehicle.
Evaluation questions, test personas and capsules from test drives never appear in prompts, exemplars or training data.

% ---------------------------------------------------------------------
\subsection{Narrators Compared}
\label{sec:conditions}

\begin{table}[t]
\caption{Narrators compared. All share Layers 1--2 and the ledger.}
\label{tab:conditions}
\begin{tabular}{lp{0.62\linewidth}}
\toprule
Code & Description \\
\midrule
T & Template NLG pipeline: rule-based content selection, aggregation, referring expressions, ISO-class lexical choice, history comparison, a fallback for every missing field (\TBD{} rules). \\
T+V & T with sampled lexical and ordering variation (controls for variety alone). \\
T$\rightarrow$P & T followed by SLM paraphrase, as in deployed driver-feedback systems~\cite{thompson2026localising}. \\
S-free & SLM given the capsule, free generation, no plan or verifier. \\
S-PGR & Plan--ground--realize with verification and template fallback, on device (ours). \\
C-PGR & S-PGR protocol with a cloud model (\TBD{}), as an upper reference. \\
\bottomrule
\end{tabular}
\end{table}

% ---------------------------------------------------------------------
\subsection{Metrics}
\label{sec:metrics}

\paragraph{Claims.}
A narrative $N$ is decomposed into atomic claims $\mathcal X(N)$ of type severity, numeric, event, anchor, comparison or hedge, in the spirit of~\cite{min2023factscore,wiseman2017challenges}.
The same LLM-based extractor is applied to every narrator.
On \TBD{} claims stratified by narrator it agreed with two annotators at $\kappa=\TBD{}$.
A claim is supported by reference $R$ when its value, unit, place and direction agree with $R$ within type-specific tolerances.
We use two references: the ledger $L$, which isolates the language layer, and ground truth $G$, which measures the system end to end:
\begin{equation}
P_R(N)=\frac{1}{|\mathcal X(N)|}\sum_{c\in\mathcal X(N)}\mathrm{supp}_R(c),\qquad R\in\{L,G\}.
\end{equation}
We report claims about facts marked \textsc{missing} separately; their rate should be zero.

\paragraph{Kinetic--semantic fidelity.}
Let $\mathcal S$ be the severity claims of a narrator pooled over the corpus.
Each claim $c\in\mathcal S$ pairs a measured intensity $x_c$ ($a_v$ for how the ride felt, $\tilde a_{w,z}$ for the road) with an expressed class $\hat k_c\in\{0,\dots,5\}$ from a fixed lexicon of ISO comfort descriptors.
With $R_c$ and $Q_c$ the mid-ranks of $x_c$ and $\hat k_c$,
\begin{equation}
\rho_{\mathrm{KSF}}=\frac{\sum_{c}(R_c-\bar R)(Q_c-\bar Q)}{\sqrt{\sum_{c}(R_c-\bar R)^2\sum_{c}(Q_c-\bar Q)^2}} .
\end{equation}
Because $\rho_{\mathrm{KSF}}$ ignores monotone distortion, we also report band consistency and an overstatement index,
\begin{equation}
\mathrm{BC}=\frac1{|\mathcal S|}\sum_{c\in\mathcal S}\ind\bigl[x_c\in\mathcal B(\hat k_c)\bigr],\qquad
\mathrm{OI}=\frac1{|\mathcal S|}\sum_{c\in\mathcal S}\bigl(\hat k_c-k^\star(x_c)\bigr),\qquad
k^\star(x)=\sum_{j=1}^5\ind[x\ge\theta_j],
\end{equation}
where $\mathcal B(k)$ is the (overlapping) ISO band of class $k$ and $\boldsymbol\theta$ are the cut-points of Section~\ref{sec:layer1}.

\paragraph{Anomaly grounding.}
Narrated event mentions $\hat{\mathcal E}$ are expanded by count; vague quantifiers expand to the lower bound of a fixed lexicon.
Each mention has a type and a resolved along-track interval $I(\hat e)$, and each reference event $e\in\mathcal E^R$ has a type and a position $s(e)$.
A mention and a reference event may match when their types agree and $\mathrm{dist}(s(e),I(\hat e))\le\delta_s$.
With $M^\star$ a maximum bipartite matching,
\begin{equation}
\mathrm{AGP}_R=\frac{|M^\star|}{|\hat{\mathcal E}|},\qquad
\mathrm{AGR}_R=\frac{|M^\star|}{|\mathcal E^R|},\qquad
\mathrm{wAGR}_R=\frac{\sum_{e\in M^\star}\mathrm{MTVV}_e}{\sum_{e\in\mathcal E^R}\mathrm{MTVV}_e}.
\end{equation}
For $R=L$, recall is computed over the must-mention events.
For $R=G$, the references are the labelled bumps (PVS) and the annotated events (phone corpus).
We use $\delta_s=25$~m, which covers label timing uncertainty and map-match error, and report 10 and 50~m in the supplement.
Crossing the two references attributes every error to perception (missed or false detection), selection (omitted) or verbalization (unsupported).

\paragraph{Calibration of hedges.}
For event claims with hedge $h\in\{\text{assertive},\text{likely},\text{possibly}\}$, let $P_G(h)$ be their precision against ground truth and $p_h$ the probability participants associate with the hedge, elicited in a short survey (cf.~\cite{kent1964wep,mielke2022linguistic}). Then
\begin{equation}
\mathrm{ECE}_{\mathrm{ling}}=\sum_h\frac{n_h}{n}\,\bigl|P_G(h)-p_h\bigr| .
\end{equation}

\paragraph{Novelty--change alignment.}
For consecutive drives $(d,d')$ of one route by one driver, let $\Delta_{\mathrm{kin}}$ be the mean absolute change of $\log a_v$ over 100~m bins plus the event-set Jaccard distance.
Let $\Delta_{\mathrm{claim}}$ be the Jaccard distance between canonicalized claim sets.
Then $\mathrm{NCA}=\mathrm{Spearman}(\Delta_{\mathrm{kin}},\Delta_{\mathrm{claim}})$ over pairs.
Lexical variation independent of content, $\mathrm{LV}=1-\mathrm{BLEU}$ between narratives with numbers and names masked, is reported alongside; variety that does not track change raises LV but not NCA.

\paragraph{Question answering and cost.}
For questions we report the correct-answer rate (judged against the ledger by raters blind to the narrator), the unsupported-claim rate, and abstention precision and recall on unanswerable questions.
System cost comprises time to first token, decode rate, end-to-end latency (p50/p95), peak memory, energy per narrative, plan validity, verifier rejection rate and template fallback rate.

% ---------------------------------------------------------------------
\subsection{Layer 1 and Layer 2}
\label{sec:l1eval}

\paragraph{Orientation.}
On PVS, the heading estimate of Eq.~\eqref{eq:yaw} recovered the known mount: $87.9^\circ$--$92.5^\circ$ for two vehicles (nominal $90^\circ$).
For the third vehicle it gave $103.1^\circ$--$105.5^\circ$ consistently across its three routes, indicating a mounting offset of about $14^\circ$.
Under 20 random three-dimensional rotations of the input per trip, the estimate rotated with the data to within $0.63^\circ$.
On the phone corpus, handling detection and re-estimation achieve \TBD{}.

\paragraph{Do maneuvers inflate vertical vibration?}
Table~\ref{tab:maneuver} compares two-second windows on asphalt (no speed bump within $\pm2$~s, 5--20~m\,s$^{-1}$) across maneuver states, controlling for speed and trip.
Braking lowered dashboard vertical vibration and turning had no detectable effect.
Acceleration added high-band energy at the dashboard that the below-suspension sensor did not see, consistent with powertrain vibration.
This is why the gate of Eq.~\eqref{eq:gate} targets acceleration and stop-and-go events and never vetoes on braking.
A wheelbase-lag cue (front then rear axle) did not separate road events: the lag between the two largest impacts did not scale with $1/v$ (Spearman $|\rho|\le0.12$, $n=60$), so it is not used.

\begin{table}[t]
\caption{Difference in log RMS vertical acceleration relative to cruising (PVS, asphalt, 2-s windows; OLS with log speed and trip effects, trip-clustered standard errors). Cruise $n=225$.}
\label{tab:maneuver}
\begin{tabular}{lcccc}
\toprule
Maneuver ($n$) & Dash 0.5--4 Hz & Dash 4--25 Hz & Wheel 4--25 Hz & Dash $-$ wheel, 4--25 Hz\\
\midrule
Brake (75) & $-0.18^{***}$ & $-0.13^{*}$ & $-0.16$ & $+0.03$\\
Accelerate (23) & $-0.18$ & $-0.06$ & $-0.37^{**}$ & $+0.31^{***}$\\
Turn (92) & $+0.07$ & $-0.04$ & $+0.06$ & $-0.10^{***}$\\
\bottomrule
\multicolumn{5}{l}{\footnotesize $^{*}p<.05$, $^{**}p<.01$, $^{***}p<.001$. Speed elasticity $\hat\gamma$: 0.32, 0.36, 0.62 (columns 1--3).}
\end{tabular}
\end{table}

\paragraph{Events and segmentation.}
Bump detection on held-out routes reaches recall \TBD{} at \TBD{} false alarms per km.
For comparison, a speed-only detector reaches \TBD{} and $\mathrm{RMS}(a_z)/\sqrt v$ reaches \TBD{}.
The gate changes false alarms by \TBD{} at fixed recall.
Against the 54 surface transitions, MI-PELT reaches boundary F1 \TBD{} at 25~m, compared with \TBD{} for constant-penalty PELT, \TBD{} for map-only boundaries and \TBD{} for PELT on speed alone.
OpenStreetMap carries a surface tag on \TBD{}\% of the matched length and a \texttt{traffic\_calming} node at \TBD{} of the 20 bumps.
Results on PVS cover three route clusters and are descriptive; inferential comparisons use the phone corpus.

% ---------------------------------------------------------------------
\subsection{Narrative Benchmark}
\label{sec:bench}

\paragraph{Tasks.}
Five task families probe where a template narrator suffices and where it does not.
\textbf{E1} answers questions from a sealed half of \TBD{} questions elicited from drivers before the system was built; the template baseline uses an intent classifier over intents designed on the other half.
\textbf{E2} writes for audiences that no template or exemplar targets (road engineer, motion-sensitive passenger, fleet manager); the template baseline uses the nearest designed audience.
\textbf{E3} narrates degraded or contradictory capsules: no gyroscope, a 60~s GNSS gap, failed map matching, map--signal disagreement, low-confidence events.
\textbf{E4} narrates repeated drives of one route.
\textbf{E5} grows the template narrator rule by rule on the development split and traces its coverage against the number of rules.
Thresholds that would count against our hypotheses were fixed before the sealed sets were used (\TBD[: registration statement]).

\paragraph{Results.}
\TBD[: Table with $P_L$, $P_G$, hallucinated-missing rate, $\rho_{\mathrm{KSF}}$, BC, OI, AGP/AGR ($L$ and $G$), $\mathrm{ECE}_{\mathrm{ling}}$, NCA, LV, QA correctness, fallback rate per narrator; frontier plot of $P_L$ against task coverage; error attribution table.]

% ---------------------------------------------------------------------
\subsection{On-Device Performance}
\label{sec:device}
\TBD[: two phones (one mid-range); model, quantization, time to first token, decode rate, end-to-end latency p50/p95, peak memory, energy per narrative, battery drain of 100--400~Hz logging, throttling over 10 consecutive generations. State which layers run on the phone.]

% ---------------------------------------------------------------------
\subsection{User Study}
\label{sec:study}

\paragraph{Design.}
Twenty-eight licensed drivers each judged narratives from four narrators (T, T+V, T$\rightarrow$P, S-PGR) in a within-subject design.
Narrator order followed a Williams design~\cite{williams1949designs} with sequences (T, T+V, S-PGR, T$\rightarrow$P), (T+V, T$\rightarrow$P, T, S-PGR), (T$\rightarrow$P, S-PGR, T+V, T) and (S-PGR, T, T$\rightarrow$P, T+V), each used by seven participants, which balances first-order carryover.
Every participant saw eight video excerpts of 2--3~min, two per narrator.
A second Latin square rotated excerpts across narrators, so each excerpt appeared with each narrator equally often and excerpts could be modelled as random effects~\cite{judd2012stimuli}.
In two of each participant's eight narratives one fact was altered (a count, a place or a severity), balanced across narrators, to measure whether errors were noticed.
Participants were debriefed about the altered narratives.

\paragraph{Participants and ethics.}
\TBD[: recruitment, inclusion criteria (licensed, $\ge$1 year of regular driving), demographics, compensation, ethics approval number]. Location traces are identifying~\cite{demontjoye2013unique}; they were stored encrypted and are released only truncated by 500~m at each end. No narrative was presented while a participant was driving.

\paragraph{Procedure.}
After consent and a practice trial, participants completed four blocks, one per narrator.
For each excerpt they watched the video, read the narrative and answered the items of Table~\ref{tab:items}.
They then answered whether anything seemed wrong, and if so, what.
Each block ended with the trust checklist of Jian et al.~\cite{jian2000trust}.
In a question-answering task, participants typed three questions about each of two further excerpts and rated the answers, which came from T with an intent classifier or from S-PGR, counterbalanced.
The session closed with a ranking of the narrators and a semi-structured interview, analysed thematically~\cite{braun2006thematic}.

\begin{table}[t]
\caption{Per-narrative items (7-point scale, 1 = strongly disagree, 7 = strongly agree).}
\label{tab:items}
\begin{tabular}{lp{0.78\linewidth}}
\toprule
Code & Item \\
\midrule
ACC1 & The summary accurately described the road conditions in the clip. \\
ACC2 & The places the summary mentioned matched where things happened in the clip. \\
SPEC & The summary was specific enough that I could find the places it mentioned again. \\
USE & This summary would help me decide how to drive this route next time. \\
CLAR & The summary was easy to understand on the first reading. \\
CERT & When the summary was unsure about something, it said so in a way that felt right. \\
NOVEL & The summary told me something about the drive that I had not noticed myself. \\
ANNOY & I would find getting a summary like this after every drive annoying. (reversed) \\
FIT & This summary was written for someone like me. \\
\bottomrule
\end{tabular}
\end{table}

\paragraph{Hypotheses.}
Four hypotheses were fixed before data collection.
H1: claim precision of S-PGR is non-inferior to T, $P_L(\text{S-PGR})-P_L(\text{T})>-0.02$.
H2: S-PGR answers the sealed questions correctly more often than T with an intent classifier.
H3: S-PGR is rated more useful (USE) than T+V.
H4: perceived accuracy, the mean of ACC1 and ACC2, of S-PGR is non-inferior to T within 0.5 scale points.
All other comparisons are exploratory.

\paragraph{Analysis.}
Ratings were analysed with cumulative link mixed models~\cite{christensen_ordinal},
\begin{equation}
\operatorname{logit}P(Y_{ij}\le k)=\theta_k-\bigl(\beta_{\mathrm{cond}(ij)}+\beta_{\mathrm{period}(ij)}+u_i+w_j\bigr),\qquad u_i\sim\mathcal N(0,\sigma_u^2),\ w_j\sim\mathcal N(0,\sigma_w^2),
\end{equation}
for participant $i$ and excerpt $j$, with Holm correction across H2 and H3 and a first-order carryover term as a check.
We also ran aligned-rank-transform ANOVA~\cite{wobbrock2011art} and Wilcoxon signed-rank tests on participant means as robustness checks.
H4 was tested on participant means of unaltered trials and holds if the lower bound of the 90\% confidence interval of the paired difference exceeds $-0.5$.
H1 used a cluster bootstrap over drives.
H2 and error detection used logistic mixed models with random intercepts for question or participant and for excerpt.

\paragraph{Power.}
With 28 participants, a paired comparison has 80\% power for standardized effects $d_z\ge0.55$ at $\alpha=.05$ ($d_z\ge0.65$ at the Holm-adjusted $\alpha=.05/3$).
The non-inferiority test of H4 has 82.5\% power when the margin equals half the standard deviation of paired differences.
The pilot ($n=\TBD{}$) estimated that standard deviation at \TBD{} scale points.

\paragraph{Field deployment.}
\TBD{} participants used RoadScribe on their own phones for 14 days.
Narrators alternated by day between T+V and S-PGR on a randomized balanced schedule.
Narratives were delivered only after the vehicle had been stationary for two minutes.
We logged opens, time to open, dwell time and questions asked, and collected three daily ratings (USE, ANNOY, NOVEL).
We model habituation as the narrator-by-day interaction in a mixed model with random slopes; with this sample size the analysis is descriptive.
 after sec3_system (shares its macros).
% Bibliography keys used below (add to your .bib):
%   menegazzo2021pvs, min2023factscore, wiseman2017challenges, kent1964wep,
%   mielke2022linguistic, williams1949designs, judd2012stimuli, jian2000trust,
%   christensen_ordinal, wobbrock2011art, braun2006thematic, demontjoye2013unique,
%   thompson2026localising, vandeemter2005templates
% =====================================================================
\providecommand{\R}{\mathbb{R}}
\providecommand{\ind}{\mathbf{1}}
\providecommand{\TBD}[1][]{\textcolor{red}{[TBD#1]}}

\section{Evaluation}
\label{sec:eval}

We evaluate RoadScribe on five questions.
\begin{description}
\item[RQ1] Does Layer~1 recover vehicle-frame vibration and road events from a phone, and are its outputs grounded in vibration rather than in speed?
\item[RQ2] Does map-informed segmentation recover road-surface boundaries better than signal-only or map-only segmentation?
\item[RQ3] How faithful are the generated narratives to the ledger and to ground truth, and how do errors split between sensing and language?
\item[RQ4] Where does a grounded SLM add coverage over a template narrator, and at what fidelity cost?
\item[RQ5] How do people judge the narratives for accuracy, usefulness, trust and annoyance, and do they notice errors?
\end{description}

% ---------------------------------------------------------------------
\subsection{Data}
\label{sec:data}

\paragraph{The PVS dataset.}
The Passive Vehicular Sensors dataset~\cite{menegazzo2021pvs} contains nine trips that cover three routes (13.8, 11.6 and 10.7~km), each driven once by three vehicle--driver pairs.
In total there are 180.1~min and 108.6~km of driving.
IMU modules at the dashboard, above the suspension and below the suspension log at 100~Hz on each side of the vehicle, GNSS logs at 1~Hz, and an environment camera records each trip.
Labels mark the surface type (asphalt, cobblestone, dirt; 54 transitions) and speed bumps (60 traversals of 20 physical bumps).
We use the left-side logger: the dashboard unit serves as the phone proxy, and the suspension units serve as reference sensors that a phone does not have.
The left and right loggers are not synchronized.
Cross-correlating their yaw rates shows offsets of up to 0.54~s that drift within a trip, so we do not average the two sides.
PVS uses fixed modules, which exercise neither phone handling nor holder dynamics.

\paragraph{Phone corpus.}
To cover what PVS cannot, we recorded \TBD{} drives (\TBD{}~km, \TBD{} vehicles) with two phones.
One phone was in a dashboard holder and recorded video; the other was fixed to the seat rail as a floor reference.
Both logged at \TBD{}~Hz and were synchronized by tap events and yaw-rate cross-correlation (residual offset \TBD{}~ms).
The corpus includes \TBD{} repetitions of one route by one driver and \TBD{} drives with deliberate phone handling.
Two annotators marked speed breakers, potholes, rough patches, surface changes and rail crossings from video without seeing any narrative (Cohen's $\kappa=\TBD{}$ on \TBD{}\% double-coded).

\paragraph{Splits.}
All learned components, thresholds and prompts are fit with routes held out.
On PVS we use leave-one-route-out folds, and for vehicle generalization we test each (route, vehicle) trip using only the four trips that share neither its route nor its vehicle.
Evaluation questions, test personas and capsules from test drives never appear in prompts, exemplars or training data.

% ---------------------------------------------------------------------
\subsection{Narrators Compared}
\label{sec:conditions}

\begin{table}[t]
\caption{Narrators compared. All share Layers 1--2 and the ledger.}
\label{tab:conditions}
\begin{tabular}{lp{0.62\linewidth}}
\toprule
Code & Description \\
\midrule
T & Template NLG pipeline: rule-based content selection, aggregation, referring expressions, ISO-class lexical choice, history comparison, a fallback for every missing field (\TBD{} rules). \\
T+V & T with sampled lexical and ordering variation (controls for variety alone). \\
T$\rightarrow$P & T followed by SLM paraphrase, as in deployed driver-feedback systems~\cite{thompson2026localising}. \\
S-free & SLM given the capsule, free generation, no plan or verifier. \\
S-PGR & Plan--ground--realize with verification and template fallback, on device (ours). \\
C-PGR & S-PGR protocol with a cloud model (\TBD{}), as an upper reference. \\
\bottomrule
\end{tabular}
\end{table}

% ---------------------------------------------------------------------
\subsection{Metrics}
\label{sec:metrics}

\paragraph{Claims.}
A narrative $N$ is decomposed into atomic claims $\mathcal X(N)$ of type severity, numeric, event, anchor, comparison or hedge, in the spirit of~\cite{min2023factscore,wiseman2017challenges}.
The same LLM-based extractor is applied to every narrator.
On \TBD{} claims stratified by narrator it agreed with two annotators at $\kappa=\TBD{}$.
A claim is supported by reference $R$ when its value, unit, place and direction agree with $R$ within type-specific tolerances.
We use two references: the ledger $L$, which isolates the language layer, and ground truth $G$, which measures the system end to end:
\begin{equation}
P_R(N)=\frac{1}{|\mathcal X(N)|}\sum_{c\in\mathcal X(N)}\mathrm{supp}_R(c),\qquad R\in\{L,G\}.
\end{equation}
We report claims about facts marked \textsc{missing} separately; their rate should be zero.

\paragraph{Kinetic--semantic fidelity.}
Let $\mathcal S$ be the severity claims of a narrator pooled over the corpus.
Each claim $c\in\mathcal S$ pairs a measured intensity $x_c$ ($a_v$ for how the ride felt, $\tilde a_{w,z}$ for the road) with an expressed class $\hat k_c\in\{0,\dots,5\}$ from a fixed lexicon of ISO comfort descriptors.
With $R_c$ and $Q_c$ the mid-ranks of $x_c$ and $\hat k_c$,
\begin{equation}
\rho_{\mathrm{KSF}}=\frac{\sum_{c}(R_c-\bar R)(Q_c-\bar Q)}{\sqrt{\sum_{c}(R_c-\bar R)^2\sum_{c}(Q_c-\bar Q)^2}} .
\end{equation}
Because $\rho_{\mathrm{KSF}}$ ignores monotone distortion, we also report band consistency and an overstatement index,
\begin{equation}
\mathrm{BC}=\frac1{|\mathcal S|}\sum_{c\in\mathcal S}\ind\bigl[x_c\in\mathcal B(\hat k_c)\bigr],\qquad
\mathrm{OI}=\frac1{|\mathcal S|}\sum_{c\in\mathcal S}\bigl(\hat k_c-k^\star(x_c)\bigr),\qquad
k^\star(x)=\sum_{j=1}^5\ind[x\ge\theta_j],
\end{equation}
where $\mathcal B(k)$ is the (overlapping) ISO band of class $k$ and $\boldsymbol\theta$ are the cut-points of Section~\ref{sec:layer1}.

\paragraph{Anomaly grounding.}
Narrated event mentions $\hat{\mathcal E}$ are expanded by count; vague quantifiers expand to the lower bound of a fixed lexicon.
Each mention has a type and a resolved along-track interval $I(\hat e)$, and each reference event $e\in\mathcal E^R$ has a type and a position $s(e)$.
A mention and a reference event may match when their types agree and $\mathrm{dist}(s(e),I(\hat e))\le\delta_s$.
With $M^\star$ a maximum bipartite matching,
\begin{equation}
\mathrm{AGP}_R=\frac{|M^\star|}{|\hat{\mathcal E}|},\qquad
\mathrm{AGR}_R=\frac{|M^\star|}{|\mathcal E^R|},\qquad
\mathrm{wAGR}_R=\frac{\sum_{e\in M^\star}\mathrm{MTVV}_e}{\sum_{e\in\mathcal E^R}\mathrm{MTVV}_e}.
\end{equation}
For $R=L$, recall is computed over the must-mention events.
For $R=G$, the references are the labelled bumps (PVS) and the annotated events (phone corpus).
We use $\delta_s=25$~m, which covers label timing uncertainty and map-match error, and report 10 and 50~m in the supplement.
Crossing the two references attributes every error to perception (missed or false detection), selection (omitted) or verbalization (unsupported).

\paragraph{Calibration of hedges.}
For event claims with hedge $h\in\{\text{assertive},\text{likely},\text{possibly}\}$, let $P_G(h)$ be their precision against ground truth and $p_h$ the probability participants associate with the hedge, elicited in a short survey (cf.~\cite{kent1964wep,mielke2022linguistic}). Then
\begin{equation}
\mathrm{ECE}_{\mathrm{ling}}=\sum_h\frac{n_h}{n}\,\bigl|P_G(h)-p_h\bigr| .
\end{equation}

\paragraph{Novelty--change alignment.}
For consecutive drives $(d,d')$ of one route by one driver, let $\Delta_{\mathrm{kin}}$ be the mean absolute change of $\log a_v$ over 100~m bins plus the event-set Jaccard distance.
Let $\Delta_{\mathrm{claim}}$ be the Jaccard distance between canonicalized claim sets.
Then $\mathrm{NCA}=\mathrm{Spearman}(\Delta_{\mathrm{kin}},\Delta_{\mathrm{claim}})$ over pairs.
Lexical variation independent of content, $\mathrm{LV}=1-\mathrm{BLEU}$ between narratives with numbers and names masked, is reported alongside; variety that does not track change raises LV but not NCA.

\paragraph{Question answering and cost.}
For questions we report the correct-answer rate (judged against the ledger by raters blind to the narrator), the unsupported-claim rate, and abstention precision and recall on unanswerable questions.
System cost comprises time to first token, decode rate, end-to-end latency (p50/p95), peak memory, energy per narrative, plan validity, verifier rejection rate and template fallback rate.

% ---------------------------------------------------------------------
\subsection{Layer 1 and Layer 2}
\label{sec:l1eval}

\paragraph{Orientation.}
On PVS, the heading estimate of Eq.~\eqref{eq:yaw} recovered the known mount: $87.9^\circ$--$92.5^\circ$ for two vehicles (nominal $90^\circ$).
For the third vehicle it gave $103.1^\circ$--$105.5^\circ$ consistently across its three routes, indicating a mounting offset of about $14^\circ$.
Under 20 random three-dimensional rotations of the input per trip, the estimate rotated with the data to within $0.63^\circ$.
On the phone corpus, handling detection and re-estimation achieve \TBD{}.

\paragraph{Do maneuvers inflate vertical vibration?}
Table~\ref{tab:maneuver} compares two-second windows on asphalt (no speed bump within $\pm2$~s, 5--20~m\,s$^{-1}$) across maneuver states, controlling for speed and trip.
Braking lowered dashboard vertical vibration and turning had no detectable effect.
Acceleration added high-band energy at the dashboard that the below-suspension sensor did not see, consistent with powertrain vibration.
This is why the gate of Eq.~\eqref{eq:gate} targets acceleration and stop-and-go events and never vetoes on braking.
A wheelbase-lag cue (front then rear axle) did not separate road events: the lag between the two largest impacts did not scale with $1/v$ (Spearman $|\rho|\le0.12$, $n=60$), so it is not used.

\begin{table}[t]
\caption{Difference in log RMS vertical acceleration relative to cruising (PVS, asphalt, 2-s windows; OLS with log speed and trip effects, trip-clustered standard errors). Cruise $n=225$.}
\label{tab:maneuver}
\begin{tabular}{lcccc}
\toprule
Maneuver ($n$) & Dash 0.5--4 Hz & Dash 4--25 Hz & Wheel 4--25 Hz & Dash $-$ wheel, 4--25 Hz\\
\midrule
Brake (75) & $-0.18^{***}$ & $-0.13^{*}$ & $-0.16$ & $+0.03$\\
Accelerate (23) & $-0.18$ & $-0.06$ & $-0.37^{**}$ & $+0.31^{***}$\\
Turn (92) & $+0.07$ & $-0.04$ & $+0.06$ & $-0.10^{***}$\\
\bottomrule
\multicolumn{5}{l}{\footnotesize $^{*}p<.05$, $^{**}p<.01$, $^{***}p<.001$. Speed elasticity $\hat\gamma$: 0.32, 0.36, 0.62 (columns 1--3).}
\end{tabular}
\end{table}

\paragraph{Events and segmentation.}
Bump detection on held-out routes reaches recall \TBD{} at \TBD{} false alarms per km.
For comparison, a speed-only detector reaches \TBD{} and $\mathrm{RMS}(a_z)/\sqrt v$ reaches \TBD{}.
The gate changes false alarms by \TBD{} at fixed recall.
Against the 54 surface transitions, MI-PELT reaches boundary F1 \TBD{} at 25~m, compared with \TBD{} for constant-penalty PELT, \TBD{} for map-only boundaries and \TBD{} for PELT on speed alone.
OpenStreetMap carries a surface tag on \TBD{}\% of the matched length and a \texttt{traffic\_calming} node at \TBD{} of the 20 bumps.
Results on PVS cover three route clusters and are descriptive; inferential comparisons use the phone corpus.

% ---------------------------------------------------------------------
\subsection{Narrative Benchmark}
\label{sec:bench}

\paragraph{Tasks.}
Five task families probe where a template narrator suffices and where it does not.
\textbf{E1} answers questions from a sealed half of \TBD{} questions elicited from drivers before the system was built; the template baseline uses an intent classifier over intents designed on the other half.
\textbf{E2} writes for audiences that no template or exemplar targets (road engineer, motion-sensitive passenger, fleet manager); the template baseline uses the nearest designed audience.
\textbf{E3} narrates degraded or contradictory capsules: no gyroscope, a 60~s GNSS gap, failed map matching, map--signal disagreement, low-confidence events.
\textbf{E4} narrates repeated drives of one route.
\textbf{E5} grows the template narrator rule by rule on the development split and traces its coverage against the number of rules.
Thresholds that would count against our hypotheses were fixed before the sealed sets were used (\TBD[: registration statement]).

\paragraph{Results.}
\TBD[: Table with $P_L$, $P_G$, hallucinated-missing rate, $\rho_{\mathrm{KSF}}$, BC, OI, AGP/AGR ($L$ and $G$), $\mathrm{ECE}_{\mathrm{ling}}$, NCA, LV, QA correctness, fallback rate per narrator; frontier plot of $P_L$ against task coverage; error attribution table.]

% ---------------------------------------------------------------------
\subsection{On-Device Performance}
\label{sec:device}
\TBD[: two phones (one mid-range); model, quantization, time to first token, decode rate, end-to-end latency p50/p95, peak memory, energy per narrative, battery drain of 100--400~Hz logging, throttling over 10 consecutive generations. State which layers run on the phone.]

% ---------------------------------------------------------------------
\subsection{User Study}
\label{sec:study}

\paragraph{Design.}
Twenty-eight licensed drivers each judged narratives from four narrators (T, T+V, T$\rightarrow$P, S-PGR) in a within-subject design.
Narrator order followed a Williams design~\cite{williams1949designs} with sequences (T, T+V, S-PGR, T$\rightarrow$P), (T+V, T$\rightarrow$P, T, S-PGR), (T$\rightarrow$P, S-PGR, T+V, T) and (S-PGR, T, T$\rightarrow$P, T+V), each used by seven participants, which balances first-order carryover.
Every participant saw eight video excerpts of 2--3~min, two per narrator.
A second Latin square rotated excerpts across narrators, so each excerpt appeared with each narrator equally often and excerpts could be modelled as random effects~\cite{judd2012stimuli}.
In two of each participant's eight narratives one fact was altered (a count, a place or a severity), balanced across narrators, to measure whether errors were noticed.
Participants were debriefed about the altered narratives.

\paragraph{Participants and ethics.}
\TBD[: recruitment, inclusion criteria (licensed, $\ge$1 year of regular driving), demographics, compensation, ethics approval number]. Location traces are identifying~\cite{demontjoye2013unique}; they were stored encrypted and are released only truncated by 500~m at each end. No narrative was presented while a participant was driving.

\paragraph{Procedure.}
After consent and a practice trial, participants completed four blocks, one per narrator.
For each excerpt they watched the video, read the narrative and answered the items of Table~\ref{tab:items}.
They then answered whether anything seemed wrong, and if so, what.
Each block ended with the trust checklist of Jian et al.~\cite{jian2000trust}.
In a question-answering task, participants typed three questions about each of two further excerpts and rated the answers, which came from T with an intent classifier or from S-PGR, counterbalanced.
The session closed with a ranking of the narrators and a semi-structured interview, analysed thematically~\cite{braun2006thematic}.

\begin{table}[t]
\caption{Per-narrative items (7-point scale, 1 = strongly disagree, 7 = strongly agree).}
\label{tab:items}
\begin{tabular}{lp{0.78\linewidth}}
\toprule
Code & Item \\
\midrule
ACC1 & The summary accurately described the road conditions in the clip. \\
ACC2 & The places the summary mentioned matched where things happened in the clip. \\
SPEC & The summary was specific enough that I could find the places it mentioned again. \\
USE & This summary would help me decide how to drive this route next time. \\
CLAR & The summary was easy to understand on the first reading. \\
CERT & When the summary was unsure about something, it said so in a way that felt right. \\
NOVEL & The summary told me something about the drive that I had not noticed myself. \\
ANNOY & I would find getting a summary like this after every drive annoying. (reversed) \\
FIT & This summary was written for someone like me. \\
\bottomrule
\end{tabular}
\end{table}

\paragraph{Hypotheses.}
Four hypotheses were fixed before data collection.
H1: claim precision of S-PGR is non-inferior to T, $P_L(\text{S-PGR})-P_L(\text{T})>-0.02$.
H2: S-PGR answers the sealed questions correctly more often than T with an intent classifier.
H3: S-PGR is rated more useful (USE) than T+V.
H4: perceived accuracy, the mean of ACC1 and ACC2, of S-PGR is non-inferior to T within 0.5 scale points.
All other comparisons are exploratory.

\paragraph{Analysis.}
Ratings were analysed with cumulative link mixed models~\cite{christensen_ordinal},
\begin{equation}
\operatorname{logit}P(Y_{ij}\le k)=\theta_k-\bigl(\beta_{\mathrm{cond}(ij)}+\beta_{\mathrm{period}(ij)}+u_i+w_j\bigr),\qquad u_i\sim\mathcal N(0,\sigma_u^2),\ w_j\sim\mathcal N(0,\sigma_w^2),
\end{equation}
for participant $i$ and excerpt $j$, with Holm correction across H2 and H3 and a first-order carryover term as a check.
We also ran aligned-rank-transform ANOVA~\cite{wobbrock2011art} and Wilcoxon signed-rank tests on participant means as robustness checks.
H4 was tested on participant means of unaltered trials and holds if the lower bound of the 90\% confidence interval of the paired difference exceeds $-0.5$.
H1 used a cluster bootstrap over drives.
H2 and error detection used logistic mixed models with random intercepts for question or participant and for excerpt.

\paragraph{Power.}
With 28 participants, a paired comparison has 80\% power for standardized effects $d_z\ge0.55$ at $\alpha=.05$ ($d_z\ge0.65$ at the Holm-adjusted $\alpha=.05/3$).
The non-inferiority test of H4 has 82.5\% power when the margin equals half the standard deviation of paired differences.
The pilot ($n=\TBD{}$) estimated that standard deviation at \TBD{} scale points.

\paragraph{Field deployment.}
\TBD{} participants used RoadScribe on their own phones for 14 days.
Narrators alternated by day between T+V and S-PGR on a randomized balanced schedule.
Narratives were delivered only after the vehicle had been stationary for two minutes.
We logged opens, time to open, dwell time and questions asked, and collected three daily ratings (USE, ANNOY, NOVEL).
We model habituation as the narrator-by-day interaction in a mixed model with random slopes; with this sample size the analysis is descriptive.
 after sec3_system (shares its macros).
% Bibliography keys used below (add to your .bib):
%   menegazzo2021pvs, min2023factscore, wiseman2017challenges, kent1964wep,
%   mielke2022linguistic, williams1949designs, judd2012stimuli, jian2000trust,
%   christensen_ordinal, wobbrock2011art, braun2006thematic, demontjoye2013unique,
%   thompson2026localising, vandeemter2005templates
% =====================================================================
\providecommand{\R}{\mathbb{R}}
\providecommand{\ind}{\mathbf{1}}
\providecommand{\TBD}[1][]{\textcolor{red}{[TBD#1]}}

\section{Evaluation}
\label{sec:eval}

We evaluate RoadScribe on five questions.
\begin{description}
\item[RQ1] Does Layer~1 recover vehicle-frame vibration and road events from a phone, and are its outputs grounded in vibration rather than in speed?
\item[RQ2] Does map-informed segmentation recover road-surface boundaries better than signal-only or map-only segmentation?
\item[RQ3] How faithful are the generated narratives to the ledger and to ground truth, and how do errors split between sensing and language?
\item[RQ4] Where does a grounded SLM add coverage over a template narrator, and at what fidelity cost?
\item[RQ5] How do people judge the narratives for accuracy, usefulness, trust and annoyance, and do they notice errors?
\end{description}

% ---------------------------------------------------------------------
\subsection{Data}
\label{sec:data}

\paragraph{The PVS dataset.}
The Passive Vehicular Sensors dataset~\cite{menegazzo2021pvs} contains nine trips that cover three routes (13.8, 11.6 and 10.7~km), each driven once by three vehicle--driver pairs.
In total there are 180.1~min and 108.6~km of driving.
IMU modules at the dashboard, above the suspension and below the suspension log at 100~Hz on each side of the vehicle, GNSS logs at 1~Hz, and an environment camera records each trip.
Labels mark the surface type (asphalt, cobblestone, dirt; 54 transitions) and speed bumps (60 traversals of 20 physical bumps).
We use the left-side logger: the dashboard unit serves as the phone proxy, and the suspension units serve as reference sensors that a phone does not have.
The left and right loggers are not synchronized.
Cross-correlating their yaw rates shows offsets of up to 0.54~s that drift within a trip, so we do not average the two sides.
PVS uses fixed modules, which exercise neither phone handling nor holder dynamics.

\paragraph{Phone corpus.}
To cover what PVS cannot, we recorded \TBD{} drives (\TBD{}~km, \TBD{} vehicles) with two phones.
One phone was in a dashboard holder and recorded video; the other was fixed to the seat rail as a floor reference.
Both logged at \TBD{}~Hz and were synchronized by tap events and yaw-rate cross-correlation (residual offset \TBD{}~ms).
The corpus includes \TBD{} repetitions of one route by one driver and \TBD{} drives with deliberate phone handling.
Two annotators marked speed breakers, potholes, rough patches, surface changes and rail crossings from video without seeing any narrative (Cohen's $\kappa=\TBD{}$ on \TBD{}\% double-coded).

\paragraph{Splits.}
All learned components, thresholds and prompts are fit with routes held out.
On PVS we use leave-one-route-out folds, and for vehicle generalization we test each (route, vehicle) trip using only the four trips that share neither its route nor its vehicle.
Evaluation questions, test personas and capsules from test drives never appear in prompts, exemplars or training data.

% ---------------------------------------------------------------------
\subsection{Narrators Compared}
\label{sec:conditions}

\begin{table}[t]
\caption{Narrators compared. All share Layers 1--2 and the ledger.}
\label{tab:conditions}
\begin{tabular}{lp{0.62\linewidth}}
\toprule
Code & Description \\
\midrule
T & Template NLG pipeline: rule-based content selection, aggregation, referring expressions, ISO-class lexical choice, history comparison, a fallback for every missing field (\TBD{} rules). \\
T+V & T with sampled lexical and ordering variation (controls for variety alone). \\
T$\rightarrow$P & T followed by SLM paraphrase, as in deployed driver-feedback systems~\cite{thompson2026localising}. \\
S-free & SLM given the capsule, free generation, no plan or verifier. \\
S-PGR & Plan--ground--realize with verification and template fallback, on device (ours). \\
C-PGR & S-PGR protocol with a cloud model (\TBD{}), as an upper reference. \\
\bottomrule
\end{tabular}
\end{table}

% ---------------------------------------------------------------------
\subsection{Metrics}
\label{sec:metrics}

\paragraph{Claims.}
A narrative $N$ is decomposed into atomic claims $\mathcal X(N)$ of type severity, numeric, event, anchor, comparison or hedge, in the spirit of~\cite{min2023factscore,wiseman2017challenges}.
The same LLM-based extractor is applied to every narrator.
On \TBD{} claims stratified by narrator it agreed with two annotators at $\kappa=\TBD{}$.
A claim is supported by reference $R$ when its value, unit, place and direction agree with $R$ within type-specific tolerances.
We use two references: the ledger $L$, which isolates the language layer, and ground truth $G$, which measures the system end to end:
\begin{equation}
P_R(N)=\frac{1}{|\mathcal X(N)|}\sum_{c\in\mathcal X(N)}\mathrm{supp}_R(c),\qquad R\in\{L,G\}.
\end{equation}
We report claims about facts marked \textsc{missing} separately; their rate should be zero.

\paragraph{Kinetic--semantic fidelity.}
Let $\mathcal S$ be the severity claims of a narrator pooled over the corpus.
Each claim $c\in\mathcal S$ pairs a measured intensity $x_c$ ($a_v$ for how the ride felt, $\tilde a_{w,z}$ for the road) with an expressed class $\hat k_c\in\{0,\dots,5\}$ from a fixed lexicon of ISO comfort descriptors.
With $R_c$ and $Q_c$ the mid-ranks of $x_c$ and $\hat k_c$,
\begin{equation}
\rho_{\mathrm{KSF}}=\frac{\sum_{c}(R_c-\bar R)(Q_c-\bar Q)}{\sqrt{\sum_{c}(R_c-\bar R)^2\sum_{c}(Q_c-\bar Q)^2}} .
\end{equation}
Because $\rho_{\mathrm{KSF}}$ ignores monotone distortion, we also report band consistency and an overstatement index,
\begin{equation}
\mathrm{BC}=\frac1{|\mathcal S|}\sum_{c\in\mathcal S}\ind\bigl[x_c\in\mathcal B(\hat k_c)\bigr],\qquad
\mathrm{OI}=\frac1{|\mathcal S|}\sum_{c\in\mathcal S}\bigl(\hat k_c-k^\star(x_c)\bigr),\qquad
k^\star(x)=\sum_{j=1}^5\ind[x\ge\theta_j],
\end{equation}
where $\mathcal B(k)$ is the (overlapping) ISO band of class $k$ and $\boldsymbol\theta$ are the cut-points of Section~\ref{sec:layer1}.

\paragraph{Anomaly grounding.}
Narrated event mentions $\hat{\mathcal E}$ are expanded by count; vague quantifiers expand to the lower bound of a fixed lexicon.
Each mention has a type and a resolved along-track interval $I(\hat e)$, and each reference event $e\in\mathcal E^R$ has a type and a position $s(e)$.
A mention and a reference event may match when their types agree and $\mathrm{dist}(s(e),I(\hat e))\le\delta_s$.
With $M^\star$ a maximum bipartite matching,
\begin{equation}
\mathrm{AGP}_R=\frac{|M^\star|}{|\hat{\mathcal E}|},\qquad
\mathrm{AGR}_R=\frac{|M^\star|}{|\mathcal E^R|},\qquad
\mathrm{wAGR}_R=\frac{\sum_{e\in M^\star}\mathrm{MTVV}_e}{\sum_{e\in\mathcal E^R}\mathrm{MTVV}_e}.
\end{equation}
For $R=L$, recall is computed over the must-mention events.
For $R=G$, the references are the labelled bumps (PVS) and the annotated events (phone corpus).
We use $\delta_s=25$~m, which covers label timing uncertainty and map-match error, and report 10 and 50~m in the supplement.
Crossing the two references attributes every error to perception (missed or false detection), selection (omitted) or verbalization (unsupported).

\paragraph{Calibration of hedges.}
For event claims with hedge $h\in\{\text{assertive},\text{likely},\text{possibly}\}$, let $P_G(h)$ be their precision against ground truth and $p_h$ the probability participants associate with the hedge, elicited in a short survey (cf.~\cite{kent1964wep,mielke2022linguistic}). Then
\begin{equation}
\mathrm{ECE}_{\mathrm{ling}}=\sum_h\frac{n_h}{n}\,\bigl|P_G(h)-p_h\bigr| .
\end{equation}

\paragraph{Novelty--change alignment.}
For consecutive drives $(d,d')$ of one route by one driver, let $\Delta_{\mathrm{kin}}$ be the mean absolute change of $\log a_v$ over 100~m bins plus the event-set Jaccard distance.
Let $\Delta_{\mathrm{claim}}$ be the Jaccard distance between canonicalized claim sets.
Then $\mathrm{NCA}=\mathrm{Spearman}(\Delta_{\mathrm{kin}},\Delta_{\mathrm{claim}})$ over pairs.
Lexical variation independent of content, $\mathrm{LV}=1-\mathrm{BLEU}$ between narratives with numbers and names masked, is reported alongside; variety that does not track change raises LV but not NCA.

\paragraph{Question answering and cost.}
For questions we report the correct-answer rate (judged against the ledger by raters blind to the narrator), the unsupported-claim rate, and abstention precision and recall on unanswerable questions.
System cost comprises time to first token, decode rate, end-to-end latency (p50/p95), peak memory, energy per narrative, plan validity, verifier rejection rate and template fallback rate.

% ---------------------------------------------------------------------
\subsection{Layer 1 and Layer 2}
\label{sec:l1eval}

\paragraph{Orientation.}
On PVS, the heading estimate of Eq.~\eqref{eq:yaw} recovered the known mount: $87.9^\circ$--$92.5^\circ$ for two vehicles (nominal $90^\circ$).
For the third vehicle it gave $103.1^\circ$--$105.5^\circ$ consistently across its three routes, indicating a mounting offset of about $14^\circ$.
Under 20 random three-dimensional rotations of the input per trip, the estimate rotated with the data to within $0.63^\circ$.
On the phone corpus, handling detection and re-estimation achieve \TBD{}.

\paragraph{Do maneuvers inflate vertical vibration?}
Table~\ref{tab:maneuver} compares two-second windows on asphalt (no speed bump within $\pm2$~s, 5--20~m\,s$^{-1}$) across maneuver states, controlling for speed and trip.
Braking lowered dashboard vertical vibration and turning had no detectable effect.
Acceleration added high-band energy at the dashboard that the below-suspension sensor did not see, consistent with powertrain vibration.
This is why the gate of Eq.~\eqref{eq:gate} targets acceleration and stop-and-go events and never vetoes on braking.
A wheelbase-lag cue (front then rear axle) did not separate road events: the lag between the two largest impacts did not scale with $1/v$ (Spearman $|\rho|\le0.12$, $n=60$), so it is not used.

\begin{table}[t]
\caption{Difference in log RMS vertical acceleration relative to cruising (PVS, asphalt, 2-s windows; OLS with log speed and trip effects, trip-clustered standard errors). Cruise $n=225$.}
\label{tab:maneuver}
\begin{tabular}{lcccc}
\toprule
Maneuver ($n$) & Dash 0.5--4 Hz & Dash 4--25 Hz & Wheel 4--25 Hz & Dash $-$ wheel, 4--25 Hz\\
\midrule
Brake (75) & $-0.18^{***}$ & $-0.13^{*}$ & $-0.16$ & $+0.03$\\
Accelerate (23) & $-0.18$ & $-0.06$ & $-0.37^{**}$ & $+0.31^{***}$\\
Turn (92) & $+0.07$ & $-0.04$ & $+0.06$ & $-0.10^{***}$\\
\bottomrule
\multicolumn{5}{l}{\footnotesize $^{*}p<.05$, $^{**}p<.01$, $^{***}p<.001$. Speed elasticity $\hat\gamma$: 0.32, 0.36, 0.62 (columns 1--3).}
\end{tabular}
\end{table}

\paragraph{Events and segmentation.}
Bump detection on held-out routes reaches recall \TBD{} at \TBD{} false alarms per km.
For comparison, a speed-only detector reaches \TBD{} and $\mathrm{RMS}(a_z)/\sqrt v$ reaches \TBD{}.
The gate changes false alarms by \TBD{} at fixed recall.
Against the 54 surface transitions, MI-PELT reaches boundary F1 \TBD{} at 25~m, compared with \TBD{} for constant-penalty PELT, \TBD{} for map-only boundaries and \TBD{} for PELT on speed alone.
OpenStreetMap carries a surface tag on \TBD{}\% of the matched length and a \texttt{traffic\_calming} node at \TBD{} of the 20 bumps.
Results on PVS cover three route clusters and are descriptive; inferential comparisons use the phone corpus.

% ---------------------------------------------------------------------
\subsection{Narrative Benchmark}
\label{sec:bench}

\paragraph{Tasks.}
Five task families probe where a template narrator suffices and where it does not.
\textbf{E1} answers questions from a sealed half of \TBD{} questions elicited from drivers before the system was built; the template baseline uses an intent classifier over intents designed on the other half.
\textbf{E2} writes for audiences that no template or exemplar targets (road engineer, motion-sensitive passenger, fleet manager); the template baseline uses the nearest designed audience.
\textbf{E3} narrates degraded or contradictory capsules: no gyroscope, a 60~s GNSS gap, failed map matching, map--signal disagreement, low-confidence events.
\textbf{E4} narrates repeated drives of one route.
\textbf{E5} grows the template narrator rule by rule on the development split and traces its coverage against the number of rules.
Thresholds that would count against our hypotheses were fixed before the sealed sets were used (\TBD[: registration statement]).

\paragraph{Results.}
\TBD[: Table with $P_L$, $P_G$, hallucinated-missing rate, $\rho_{\mathrm{KSF}}$, BC, OI, AGP/AGR ($L$ and $G$), $\mathrm{ECE}_{\mathrm{ling}}$, NCA, LV, QA correctness, fallback rate per narrator; frontier plot of $P_L$ against task coverage; error attribution table.]

% ---------------------------------------------------------------------
\subsection{On-Device Performance}
\label{sec:device}
\TBD[: two phones (one mid-range); model, quantization, time to first token, decode rate, end-to-end latency p50/p95, peak memory, energy per narrative, battery drain of 100--400~Hz logging, throttling over 10 consecutive generations. State which layers run on the phone.]

% ---------------------------------------------------------------------
\subsection{User Study}
\label{sec:study}

\paragraph{Design.}
Twenty-eight licensed drivers each judged narratives from four narrators (T, T+V, T$\rightarrow$P, S-PGR) in a within-subject design.
Narrator order followed a Williams design~\cite{williams1949designs} with sequences (T, T+V, S-PGR, T$\rightarrow$P), (T+V, T$\rightarrow$P, T, S-PGR), (T$\rightarrow$P, S-PGR, T+V, T) and (S-PGR, T, T$\rightarrow$P, T+V), each used by seven participants, which balances first-order carryover.
Every participant saw eight video excerpts of 2--3~min, two per narrator.
A second Latin square rotated excerpts across narrators, so each excerpt appeared with each narrator equally often and excerpts could be modelled as random effects~\cite{judd2012stimuli}.
In two of each participant's eight narratives one fact was altered (a count, a place or a severity), balanced across narrators, to measure whether errors were noticed.
Participants were debriefed about the altered narratives.

\paragraph{Participants and ethics.}
\TBD[: recruitment, inclusion criteria (licensed, $\ge$1 year of regular driving), demographics, compensation, ethics approval number]. Location traces are identifying~\cite{demontjoye2013unique}; they were stored encrypted and are released only truncated by 500~m at each end. No narrative was presented while a participant was driving.

\paragraph{Procedure.}
After consent and a practice trial, participants completed four blocks, one per narrator.
For each excerpt they watched the video, read the narrative and answered the items of Table~\ref{tab:items}.
They then answered whether anything seemed wrong, and if so, what.
Each block ended with the trust checklist of Jian et al.~\cite{jian2000trust}.
In a question-answering task, participants typed three questions about each of two further excerpts and rated the answers, which came from T with an intent classifier or from S-PGR, counterbalanced.
The session closed with a ranking of the narrators and a semi-structured interview, analysed thematically~\cite{braun2006thematic}.

\begin{table}[t]
\caption{Per-narrative items (7-point scale, 1 = strongly disagree, 7 = strongly agree).}
\label{tab:items}
\begin{tabular}{lp{0.78\linewidth}}
\toprule
Code & Item \\
\midrule
ACC1 & The summary accurately described the road conditions in the clip. \\
ACC2 & The places the summary mentioned matched where things happened in the clip. \\
SPEC & The summary was specific enough that I could find the places it mentioned again. \\
USE & This summary would help me decide how to drive this route next time. \\
CLAR & The summary was easy to understand on the first reading. \\
CERT & When the summary was unsure about something, it said so in a way that felt right. \\
NOVEL & The summary told me something about the drive that I had not noticed myself. \\
ANNOY & I would find getting a summary like this after every drive annoying. (reversed) \\
FIT & This summary was written for someone like me. \\
\bottomrule
\end{tabular}
\end{table}

\paragraph{Hypotheses.}
Four hypotheses were fixed before data collection.
H1: claim precision of S-PGR is non-inferior to T, $P_L(\text{S-PGR})-P_L(\text{T})>-0.02$.
H2: S-PGR answers the sealed questions correctly more often than T with an intent classifier.
H3: S-PGR is rated more useful (USE) than T+V.
H4: perceived accuracy, the mean of ACC1 and ACC2, of S-PGR is non-inferior to T within 0.5 scale points.
All other comparisons are exploratory.

\paragraph{Analysis.}
Ratings were analysed with cumulative link mixed models~\cite{christensen_ordinal},
\begin{equation}
\operatorname{logit}P(Y_{ij}\le k)=\theta_k-\bigl(\beta_{\mathrm{cond}(ij)}+\beta_{\mathrm{period}(ij)}+u_i+w_j\bigr),\qquad u_i\sim\mathcal N(0,\sigma_u^2),\ w_j\sim\mathcal N(0,\sigma_w^2),
\end{equation}
for participant $i$ and excerpt $j$, with Holm correction across H2 and H3 and a first-order carryover term as a check.
We also ran aligned-rank-transform ANOVA~\cite{wobbrock2011art} and Wilcoxon signed-rank tests on participant means as robustness checks.
H4 was tested on participant means of unaltered trials and holds if the lower bound of the 90\% confidence interval of the paired difference exceeds $-0.5$.
H1 used a cluster bootstrap over drives.
H2 and error detection used logistic mixed models with random intercepts for question or participant and for excerpt.

\paragraph{Power.}
With 28 participants, a paired comparison has 80\% power for standardized effects $d_z\ge0.55$ at $\alpha=.05$ ($d_z\ge0.65$ at the Holm-adjusted $\alpha=.05/3$).
The non-inferiority test of H4 has 82.5\% power when the margin equals half the standard deviation of paired differences.
The pilot ($n=\TBD{}$) estimated that standard deviation at \TBD{} scale points.

\paragraph{Field deployment.}
\TBD{} participants used RoadScribe on their own phones for 14 days.
Narrators alternated by day between T+V and S-PGR on a randomized balanced schedule.
Narratives were delivered only after the vehicle had been stationary for two minutes.
We logged opens, time to open, dwell time and questions asked, and collected three daily ratings (USE, ANNOY, NOVEL).
We model habituation as the narrator-by-day interaction in a mixed model with random slopes; with this sample size the analysis is descriptive.
